% --- Archon physics formalization source begin ---
% archon:physics
% archon:covers PhyXMiniProblems/problem_phyx_mini_0499.lean
% archon:source-report reports/phyx_mini/problem_phyx_mini_0499.source.json

\chapter{Physics problem phyx\_mini\_0499}
\label{ch:PhyXMiniProblems_problem_phyx_mini_0499}

\paragraph{Problem source.}
\#\# Physical scenario

The graph in figure measured in a photoelectric-effect experiment.

\#\# Question

What experimental value of Planck’s constant is obtained from these data?

\#\# Answer choices

- A: A: \textbackslash{}( 6.8 \textbackslash{}times 10\textasciicircum{}\{-34\} \textbackslash{}text\{ Js\} \textbackslash{})

- B: B: \textbackslash{}( 3.2 \textbackslash{}times 10\textasciicircum{}\{-34\} \textbackslash{}text\{ Js\} \textbackslash{})

- C: C: \textbackslash{}( 6.4 \textbackslash{}times 10\textasciicircum{}\{-34\} \textbackslash{}text\{ Js\} \textbackslash{})

- D: D: \textbackslash{}( 7.4 \textbackslash{}times 10\textasciicircum{}\{-34\} \textbackslash{}text\{ Js\} \textbackslash{})

\#\# Figure caption (auxiliary; use the image as primary evidence)

The image is a graph illustrating the relationship between stopping voltage (\textbackslash{}(V\_\{\textbackslash{}text\{stop\}\}\textbackslash{})) and frequency (\textbackslash{}(f\textbackslash{})). 

- **X-Axis**: Represents frequency (\textbackslash{}(f\textbackslash{})) in units of \textbackslash{}(× 10\textasciicircum{}\{15\}\textbackslash{}) Hz. It is labeled starting from 0 to 3.
- **Y-Axis**: Represents stopping voltage (\textbackslash{}(V\_\{\textbackslash{}text\{stop\}\}\textbackslash{})) in volts (V). It is labeled from 0 to 8.
- **Line**: A green line shows a positive linear relationship, starting at approximately (0.6,0) and extending to about (3,8).

The graph suggests that as the frequency increases, the stopping voltage also increases linearly.

\#\# Dataset metadata

Quantum Phenomena; Physical Model Grounding Reasoning, Numerical Reasoning

\paragraph{Recorded answer/context.}
C: C: \textbackslash{}( 6.4 \textbackslash{}times 10\textasciicircum{}\{-34\} \textbackslash{}text\{ Js\} \textbackslash{})

\paragraph{Figure/image path.}
/root/proposal\_for\_physic/science-mango/phyx\_mini\_run/phyx\_data/test\_image/499.png

\paragraph{Formalization target.}
create a compiling Lean file with sorry bodies at `PhyXMiniProblems/problem\_phyx\_mini\_0499.lean`. The Lean declarations must preserve the physical quantities, dimensions or dimensional roles, figure labels, governing-law hypotheses, and final relation expressed by this problem.
Use Mathlib/Physlib names found through LeanExplore where available. If a physics API is missing, introduce faithful local abstractions rather than scalar placeholder aliases.

\begin{theorem}[Physics formalization target]
\label{thm:physics:phyx_mini_0499:target}
\lean{PhyXMiniProblems.ProblemPhyXMini0499.problem_phyx_mini_0499}
\uses{def:physics:phyx-mini-0499:phyxminiproblems-problemphyxmini0499-actioninjouleseconds, def:physics:phyx-mini-0499:phyxminiproblems-problemphyxmini0499-photoelectricexperiment, def:physics:phyx-mini-0499:phyxminiproblems-problemphyxmini0499-matchesprimaryphotoelectricgraph, def:physics:phyx-mini-0499:phyxminiproblems-problemphyxmini0499-hasphysicalphotoelectricparameters, def:physics:phyx-mini-0499:phyxminiproblems-problemphyxmini0499-satisfieseinsteinphotoelectricequation, def:physics:phyx-mini-0499:phyxminiproblems-problemphyxmini0499-usestextbookelementarychargecalibration, def:physics:phyx-mini-0499:phyxminiproblems-problemphyxmini0499-answerplanckconstantinjouleseconds}
The assigned autoformalize agent should translate this physics problem into Lean declarations in the covered file, with theorem and lemma proof bodies written as `by sorry`.
\end{theorem}
\begin{proof}
This is an autoformalization task, not a proof task. Produce statements that can later be proved by the physics prover without weakening the physical model.
\end{proof}

% --- Archon declaration topology begin ---
\section{Declaration topology}
\begin{definition}[electric Potential Dimension]
  \label{def:physics:phyx-mini-0499:phyxminiproblems-problemphyxmini0499-electricpotentialdimension}
  \lean{PhyXMiniProblems.ProblemPhyXMini0499.electricPotentialDimension}
  The SI dimension of electric potential, energy per unit charge
\end{definition}
\begin{proof}
  This is the declared model definition of electric Potential Dimension.
\end{proof}
\begin{definition}[action Dimension]
  \label{def:physics:phyx-mini-0499:phyxminiproblems-problemphyxmini0499-actiondimension}
  \lean{PhyXMiniProblems.ProblemPhyXMini0499.actionDimension}
  The SI dimension of action, equivalently joule-seconds
\end{definition}
\begin{proof}
  This is the declared model definition of action Dimension.
\end{proof}
\begin{definition}[Frequency Quantity]
  \label{def:physics:phyx-mini-0499:phyxminiproblems-problemphyxmini0499-frequencyquantity}
  \lean{PhyXMiniProblems.ProblemPhyXMini0499.FrequencyQuantity}
  A nonnegative physical frequency carrying inverse-time dimension
\end{definition}
\begin{proof}
  This is the declared model definition of Frequency Quantity.
\end{proof}
\begin{definition}[Electric Potential Quantity]
  \label{def:physics:phyx-mini-0499:phyxminiproblems-problemphyxmini0499-electricpotentialquantity}
  \lean{PhyXMiniProblems.ProblemPhyXMini0499.ElectricPotentialQuantity}
  \uses{def:physics:phyx-mini-0499:phyxminiproblems-problemphyxmini0499-electricpotentialdimension}
  A nonnegative physical electric-potential magnitude
\end{definition}
\begin{proof}
  This is the declared model definition of Electric Potential Quantity.
\end{proof}
\begin{definition}[Charge Magnitude Quantity]
  \label{def:physics:phyx-mini-0499:phyxminiproblems-problemphyxmini0499-chargemagnitudequantity}
  \lean{PhyXMiniProblems.ProblemPhyXMini0499.ChargeMagnitudeQuantity}
  A nonnegative physical electric-charge magnitude
\end{definition}
\begin{proof}
  This is the declared model definition of Charge Magnitude Quantity.
\end{proof}
\begin{definition}[Planck Constant Quantity]
  \label{def:physics:phyx-mini-0499:phyxminiproblems-problemphyxmini0499-planckconstantquantity}
  \lean{PhyXMiniProblems.ProblemPhyXMini0499.PlanckConstantQuantity}
  \uses{def:physics:phyx-mini-0499:phyxminiproblems-problemphyxmini0499-actiondimension}
  A nonnegative physical quantity with the dimension of Planck's constant
\end{definition}
\begin{proof}
  This is the declared model definition of Planck Constant Quantity.
\end{proof}
\begin{definition}[frequency In Hertz]
  \label{def:physics:phyx-mini-0499:phyxminiproblems-problemphyxmini0499-frequencyinhertz}
  \lean{PhyXMiniProblems.ProblemPhyXMini0499.frequencyInHertz}
  \uses{def:physics:phyx-mini-0499:phyxminiproblems-problemphyxmini0499-frequencyquantity}
  Hertz readout of a physical frequency in the SI unit choice
\end{definition}
\begin{proof}
  This is the declared model definition of frequency In Hertz.
\end{proof}
\begin{definition}[electric Potential In Volts]
  \label{def:physics:phyx-mini-0499:phyxminiproblems-problemphyxmini0499-electricpotentialinvolts}
  \lean{PhyXMiniProblems.ProblemPhyXMini0499.electricPotentialInVolts}
  \uses{def:physics:phyx-mini-0499:phyxminiproblems-problemphyxmini0499-electricpotentialquantity}
  Volt readout of a physical electric potential in the SI unit choice
\end{definition}
\begin{proof}
  This is the declared model definition of electric Potential In Volts.
\end{proof}
\begin{definition}[charge Magnitude In Coulombs]
  \label{def:physics:phyx-mini-0499:phyxminiproblems-problemphyxmini0499-chargemagnitudeincoulombs}
  \lean{PhyXMiniProblems.ProblemPhyXMini0499.chargeMagnitudeInCoulombs}
  \uses{def:physics:phyx-mini-0499:phyxminiproblems-problemphyxmini0499-chargemagnitudequantity}
  Coulomb readout of a physical charge magnitude in the SI unit choice
\end{definition}
\begin{proof}
  This is the declared model definition of charge Magnitude In Coulombs.
\end{proof}
\begin{definition}[action In Joule Seconds]
  \label{def:physics:phyx-mini-0499:phyxminiproblems-problemphyxmini0499-actioninjouleseconds}
  \lean{PhyXMiniProblems.ProblemPhyXMini0499.actionInJouleSeconds}
  \uses{def:physics:phyx-mini-0499:phyxminiproblems-problemphyxmini0499-planckconstantquantity}
  Joule-second readout of a physical action in the SI unit choice
\end{definition}
\begin{proof}
  This is the declared model definition of action In Joule Seconds.
\end{proof}
\begin{definition}[energy In Joules]
  \label{def:physics:phyx-mini-0499:phyxminiproblems-problemphyxmini0499-energyinjoules}
  \lean{PhyXMiniProblems.ProblemPhyXMini0499.energyInJoules}
  Joule readout of a signed dimensionful energy
\end{definition}
\begin{proof}
  This is the declared model definition of energy In Joules.
\end{proof}
\begin{definition}[Graph Axis Quantity]
  \label{def:physics:phyx-mini-0499:phyxminiproblems-problemphyxmini0499-graphaxisquantity}
  \lean{PhyXMiniProblems.ProblemPhyXMini0499.GraphAxisQuantity}
  The two physical quantities assigned to the graph axes
\end{definition}
\begin{proof}
  This is the declared model definition of Graph Axis Quantity.
\end{proof}
\begin{definition}[Graph Display Unit]
  \label{def:physics:phyx-mini-0499:phyxminiproblems-problemphyxmini0499-graphdisplayunit}
  \lean{PhyXMiniProblems.ProblemPhyXMini0499.GraphDisplayUnit}
  The two unit labels printed beside the graph axes
\end{definition}
\begin{proof}
  This is the declared model definition of Graph Display Unit.
\end{proof}
\begin{definition}[Graph Line Color]
  \label{def:physics:phyx-mini-0499:phyxminiproblems-problemphyxmini0499-graphlinecolor}
  \lean{PhyXMiniProblems.ProblemPhyXMini0499.GraphLineColor}
  The visible color of the data line
\end{definition}
\begin{proof}
  This is the declared model definition of Graph Line Color.
\end{proof}
\begin{definition}[Graph Line Trend]
  \label{def:physics:phyx-mini-0499:phyxminiproblems-problemphyxmini0499-graphlinetrend}
  \lean{PhyXMiniProblems.ProblemPhyXMini0499.GraphLineTrend}
  Qualitative direction of the displayed line as frequency increases
\end{definition}
\begin{proof}
  This is the declared model definition of Graph Line Trend.
\end{proof}
\begin{definition}[Photoelectric Plot Point]
  \label{def:physics:phyx-mini-0499:phyxminiproblems-problemphyxmini0499-photoelectricplotpoint}
  \lean{PhyXMiniProblems.ProblemPhyXMini0499.PhotoelectricPlotPoint}
  A point expressed in the scalar coordinates printed on the graph
\end{definition}
\begin{proof}
  This is the declared model definition of Photoelectric Plot Point.
\end{proof}
\begin{definition}[Photoelectric Stopping Potential Graph]
  \label{def:physics:phyx-mini-0499:phyxminiproblems-problemphyxmini0499-photoelectricstoppingpotentialgraph}
  \lean{PhyXMiniProblems.ProblemPhyXMini0499.PhotoelectricStoppingPotentialGraph}
  \uses{def:physics:phyx-mini-0499:phyxminiproblems-problemphyxmini0499-graphaxisquantity, def:physics:phyx-mini-0499:phyxminiproblems-problemphyxmini0499-graphdisplayunit, def:physics:phyx-mini-0499:phyxminiproblems-problemphyxmini0499-graphlinecolor, def:physics:phyx-mini-0499:phyxminiproblems-problemphyxmini0499-graphlinetrend, def:physics:phyx-mini-0499:phyxminiproblems-problemphyxmini0499-photoelectricplotpoint}
  The graph records literal display coordinates separately from the physical observables of the experiment. The largest labeled ticks are 3 and 8; the green line itself extends slightly past them in the raster
\end{definition}
\begin{proof}
  This is the declared model definition of Photoelectric Stopping Potential Graph.
\end{proof}
\begin{definition}[Photoelectric Experiment]
  \label{def:physics:phyx-mini-0499:phyxminiproblems-problemphyxmini0499-photoelectricexperiment}
  \lean{PhyXMiniProblems.ProblemPhyXMini0499.PhotoelectricExperiment}
  \uses{def:physics:phyx-mini-0499:phyxminiproblems-problemphyxmini0499-frequencyquantity, def:physics:phyx-mini-0499:phyxminiproblems-problemphyxmini0499-electricpotentialquantity, def:physics:phyx-mini-0499:phyxminiproblems-problemphyxmini0499-chargemagnitudequantity, def:physics:phyx-mini-0499:phyxminiproblems-problemphyxmini0499-planckconstantquantity, def:physics:phyx-mini-0499:phyxminiproblems-problemphyxmini0499-photoelectricstoppingpotentialgraph}
  The two reference frequencies name the physical measurements corresponding to the readable graph points. The experimental Planck constant is an unknown physical parameter, while the electron charge magnitude and work function are independent physical inputs to the Einstein photoelectric equation
\end{definition}
\begin{proof}
  This is the declared model definition of Photoelectric Experiment.
\end{proof}
\begin{definition}[plotted Measurement]
  \label{def:physics:phyx-mini-0499:phyxminiproblems-problemphyxmini0499-plottedmeasurement}
  \lean{PhyXMiniProblems.ProblemPhyXMini0499.plottedMeasurement}
  \uses{def:physics:phyx-mini-0499:phyxminiproblems-problemphyxmini0499-frequencyquantity, def:physics:phyx-mini-0499:phyxminiproblems-problemphyxmini0499-frequencyinhertz, def:physics:phyx-mini-0499:phyxminiproblems-problemphyxmini0499-electricpotentialinvolts, def:physics:phyx-mini-0499:phyxminiproblems-problemphyxmini0499-photoelectricplotpoint, def:physics:phyx-mini-0499:phyxminiproblems-problemphyxmini0499-photoelectricexperiment}
  The scalar plot coordinate associated with a measured physical frequency
\end{definition}
\begin{proof}
  This is the declared model definition of plotted Measurement.
\end{proof}
\begin{definition}[Matches Primary Photoelectric Graph]
  \label{def:physics:phyx-mini-0499:phyxminiproblems-problemphyxmini0499-matchesprimaryphotoelectricgraph}
  \lean{PhyXMiniProblems.ProblemPhyXMini0499.MatchesPrimaryPhotoelectricGraph}
  \uses{def:physics:phyx-mini-0499:phyxminiproblems-problemphyxmini0499-frequencyinhertz, def:physics:phyx-mini-0499:phyxminiproblems-problemphyxmini0499-electricpotentialinvolts, def:physics:phyx-mini-0499:phyxminiproblems-problemphyxmini0499-photoelectricplotpoint, def:physics:phyx-mini-0499:phyxminiproblems-problemphyxmini0499-photoelectricexperiment, def:physics:phyx-mini-0499:phyxminiproblems-problemphyxmini0499-plottedmeasurement}
  Axis assignments, labels, qualitative line geometry, and the two clear reference points read from the primary image. The auxiliary caption's approximate horizontal intercept 0.6 is not used because the raster itself clearly places the intercept at the tick labeled 1
\end{definition}
\begin{proof}
  This is the declared model definition of Matches Primary Photoelectric Graph.
\end{proof}
\begin{definition}[Has Physical Photoelectric Parameters]
  \label{def:physics:phyx-mini-0499:phyxminiproblems-problemphyxmini0499-hasphysicalphotoelectricparameters}
  \lean{PhyXMiniProblems.ProblemPhyXMini0499.HasPhysicalPhotoelectricParameters}
  \uses{def:physics:phyx-mini-0499:phyxminiproblems-problemphyxmini0499-chargemagnitudeincoulombs, def:physics:phyx-mini-0499:phyxminiproblems-problemphyxmini0499-actioninjouleseconds, def:physics:phyx-mini-0499:phyxminiproblems-problemphyxmini0499-energyinjoules, def:physics:phyx-mini-0499:phyxminiproblems-problemphyxmini0499-photoelectricexperiment}
  Positivity conditions for the physical parameters of the experiment
\end{definition}
\begin{proof}
  This is the declared model definition of Has Physical Photoelectric Parameters.
\end{proof}
\begin{definition}[Satisfies Einstein Photoelectric Equation]
  \label{def:physics:phyx-mini-0499:phyxminiproblems-problemphyxmini0499-satisfieseinsteinphotoelectricequation}
  \lean{PhyXMiniProblems.ProblemPhyXMini0499.SatisfiesEinsteinPhotoelectricEquation}
  \uses{def:physics:phyx-mini-0499:phyxminiproblems-problemphyxmini0499-frequencyinhertz, def:physics:phyx-mini-0499:phyxminiproblems-problemphyxmini0499-electricpotentialinvolts, def:physics:phyx-mini-0499:phyxminiproblems-problemphyxmini0499-chargemagnitudeincoulombs, def:physics:phyx-mini-0499:phyxminiproblems-problemphyxmini0499-actioninjouleseconds, def:physics:phyx-mini-0499:phyxminiproblems-problemphyxmini0499-energyinjoules, def:physics:phyx-mini-0499:phyxminiproblems-problemphyxmini0499-photoelectricexperiment}
  The Einstein photoelectric equation at every measured frequency: e V\_stop = h f - phi. All terms are SI scalar readouts of dimensionally typed physical quantities, so each side represents an energy in joules. This governing law contains no numerical value for the unknown experimental Planck constant
\end{definition}
\begin{proof}
  This is the declared model definition of Satisfies Einstein Photoelectric Equation.
\end{proof}
\begin{definition}[Uses Textbook Elementary Charge Calibration]
  \label{def:physics:phyx-mini-0499:phyxminiproblems-problemphyxmini0499-usestextbookelementarychargecalibration}
  \lean{PhyXMiniProblems.ProblemPhyXMini0499.UsesTextbookElementaryChargeCalibration}
  \uses{def:physics:phyx-mini-0499:phyxminiproblems-problemphyxmini0499-chargemagnitudeincoulombs, def:physics:phyx-mini-0499:phyxminiproblems-problemphyxmini0499-photoelectricexperiment}
  The two-significant-figure elementary-charge calibration implicit in the recorded multiple-choice computation. The exact SI charge would change only digits beyond the precision supported by the plotted readouts
\end{definition}
\begin{proof}
  This is the declared model definition of Uses Textbook Elementary Charge Calibration.
\end{proof}
\begin{definition}[Answer Choice]
  \label{def:physics:phyx-mini-0499:phyxminiproblems-problemphyxmini0499-answerchoice}
  \lean{PhyXMiniProblems.ProblemPhyXMini0499.AnswerChoice}
  Labels of the four answer choices in the problem statement
\end{definition}
\begin{proof}
  This is the declared model definition of Answer Choice.
\end{proof}
\begin{definition}[answer Planck Constant In Joule Seconds]
  \label{def:physics:phyx-mini-0499:phyxminiproblems-problemphyxmini0499-answerplanckconstantinjouleseconds}
  \lean{PhyXMiniProblems.ProblemPhyXMini0499.answerPlanckConstantInJouleSeconds}
  \uses{def:physics:phyx-mini-0499:phyxminiproblems-problemphyxmini0499-answerchoice}
  Planck-constant readout printed beside each choice, in joule-seconds
\end{definition}
\begin{proof}
  This is the declared model definition of answer Planck Constant In Joule Seconds.
\end{proof}
\begin{definition}[stopping Potential Slope In Volt Seconds]
  \label{def:physics:phyx-mini-0499:phyxminiproblems-problemphyxmini0499-stoppingpotentialslopeinvoltseconds}
  \lean{PhyXMiniProblems.ProblemPhyXMini0499.stoppingPotentialSlopeInVoltSeconds}
  \uses{def:physics:phyx-mini-0499:phyxminiproblems-problemphyxmini0499-frequencyinhertz, def:physics:phyx-mini-0499:phyxminiproblems-problemphyxmini0499-electricpotentialinvolts, def:physics:phyx-mini-0499:phyxminiproblems-problemphyxmini0499-photoelectricexperiment}
  The two image readouts determine the stopping-potential-versus-frequency slope 8 V / (2 * 10\textasciicircum{}15 Hz) = 4 * 10\textasciicircum{}-15 V s
\end{definition}
\begin{proof}
  This is the declared model definition of stopping Potential Slope In Volt Seconds.
\end{proof}
\begin{lemma}[stopping Potential Slope from primary graph]
  \label{lem:physics:phyx-mini-0499:phyxminiproblems-problemphyxmini0499-stoppingpotentialslope-from-primary-graph}
  \lean{PhyXMiniProblems.ProblemPhyXMini0499.stoppingPotentialSlope_from_primary_graph}
  \uses{def:physics:phyx-mini-0499:phyxminiproblems-problemphyxmini0499-photoelectricexperiment, def:physics:phyx-mini-0499:phyxminiproblems-problemphyxmini0499-matchesprimaryphotoelectricgraph, def:physics:phyx-mini-0499:phyxminiproblems-problemphyxmini0499-stoppingpotentialslopeinvoltseconds}
  The physical slope derived solely from the two primary-image points
\end{lemma}
\begin{proof}
  The conclusion follows from the governing relations and figure readouts named in the statement of stopping Potential Slope from primary graph.
\end{proof}
% --- Archon declaration topology end ---
% --- Archon physics formalization source end ---
